% Modernised reading — fonds Alexandre Grothendieck, shelfmark 161-5
% (pages 1-28, the whole folder).
%
% This is an INTERPRETATION, not a transcription. It re-reads the pages in
% today's notation and vocabulary, states the hypotheses the manuscript leaves
% implicit, and drops the critical apparatus so that the mathematics reads
% straight through. Everything the brackets and underlines carried moves into a
% footnote. It is held to being mathematically correct as it stands; where the
% manuscript is loose or mistaken, the true statement is what appears here and
% the footnote says what the page has. In any dispute of substance the
% transcription governs: whoever cites Grothendieck must cite
% batch-01.fr.tex (pp. 1-20) or batch-02.fr.tex (pp. 21-28).
%
% Pass: Opus 5.5 (claude-opus-5-5), 2026-09-22 — first pass, unchecked against the pages by a human.
% The folder was read whole, both batches together; this is one document for
% the shelfmark. It was made on Opus 5.5 and has not been measured against the
% reference reading of folder 115 (made on Opus 5).
% Revised 2026-09-23 (Opus 5.5): (1) the K/L exchange — the account now matches
% the transcription, which records it from p. 23 c) as what the pages write
% (checked on the facsimile), and lists p. 24's « Q_L descendant Q_Omega »
% among the instances; (2) wording that ordered the leaves in time with 1970
% (« l'appellera en 1970 », « répondra », « posera », « Tate venait de »,
% « postérieur au dossier », future tenses) or inferred what the pages knew
% (« ne les connaissent pas », « n'est connu des pages », « pouvait savoir »)
% rephrased to say what the pages contain and do not cite; (3) the p. 9
% footnote: the description of M as the smallest group whose extension to C
% contains the image of h is the p. 19 note's (about H_V°), not the p. 9
% note's, and C is the completion of Kbar, not the complex numbers.
%
% Spine. pp. 1-21: two printed offprints of J.-P. Serre (Driebergen 1966, pub.
% 1967; Annuaire du Collège de France, course 1967-68), carrying five red-ink
% annotations in a hand NOT established as Grothendieck's; they cite S. Sen
% (Invent. Math., Annals) and T. Springer's Jordan-algebras book — presumably
% 1972-73 publications, which would put the annotations after 1972 (inference,
% said as such; \dating is left as Montpellier gives it). pp. 22-28: his
% manuscript. Hodge-Tate representations as a Tannakian category, Tate's
% embedding of BT groups up to isogeny, the Hodge-Tate structure written as
% torsors over (Omega, pi), the crystalline data a)-d) and the Question of p. 23
% (extend them as ⊗-functors?); pp. 24-26 axiomatise this as data 1)-6) for an
% abstract G (p. 24 abandoned, pp. 25-26 complete), ending in an action of an
% extension Gamma of Z on the de Rham torsor; pp. 27-28 translate uniqueness
% into non-abelian H^1 and ask a Question; lower p. 28, a separate ink
% computation (Néron lft model of R_{K/Q_p} G_m).
%
% Notation fixed for the whole document. K the base field, K_0 = W(k)[1/p]
% (the page's L on p. 22 and p. 23 a)-b), its K from p. 23 c) on), O_K (the
% page's V), sigma the Frobenius of K_0,
% Omega, pi = Gal(Kbar/K), pi' = Gal(Kbar/K_0). omega(M) = M ⊗ Omega with pi on
% Omega only; omega^!(M) = M ⊗ Omega with diagonal action (p. 22 writes M_Omega
% for the latter, p. 25 M^!_Omega). G^Hdg, U^Hdg for the page's G'', U'' (p. 22
% writes G^Hdg once). A = G'(Omega) for the page's G'_S(S). Uniformiser: varpi
% (the page's arrow « pi » on p. 28).
%
% Substantive departures, with pages.
%  p. 22  M^(n): the heading reads (n) but the text says degrees 0 and 1; we
%         take M^(n) = weights in [0,n], M^(1) = weights {0,1}.
%  p. 22  M^BT « engendrée par M^(1)_K »: we take the Tannakian subcategory
%         generated by the image of Tate's embedding, which is what the
%         Question of p. 23 needs; M^(1) is strictly larger (Tate curve).
%  p. 23c-28  K and L exchange roles: from T_DR = T_cr ⊗_K L (p. 23) on, the
%         page's L is the base field and its K the unramified subfield
%         (Q_L « descendant Q_Omega » p. 24, pi = Aut_L(Omega) p. 25, descent
%         « de L : K » and P_K^(sigma?) p. 26, G'_S(S)^pi = G'_L(L) pp. 27-28).
%         The transcription records the exchange as what the pages write and
%         flags it once, at p. 23 c). We write K and K_0 throughout.
%  p. 24-25  the twist by T vs T^*: fixed by the formula of p. 25
%         (gr^i ⊗ Omega(-i)).
%  p. 25  j^! pi-invariant does not by itself say Hodge-Tate; the Axiom does.
%  p. 26  G'_L ≃ G_L « non canoniquement » holds when P is a trivial torsor
%         (e.g. G = GL_n); supplied. U'_L on p. 26 read as U''.
%  p. 24  u.alpha = tau(alpha)^{-1} u tau(alpha): letters evidently exchanged.
%  p. 28  « il suffirait »: the condition is also necessary (given one
%         extension); stated as an equivalence.
%  p. 28  the Question, as posed, has a negative answer (G = G_m example, ours).
%  p. 28 (ink)  the component group is the constant Z only for K/Q_p totally
%         ramified; the map from N° x Z to N is an iso on Z_p-points and special
%         fibres, not on generic fibres.
%  p. 9 (annotation)  M = I_alg holds for the neutral component only, if M is
%         the connected group of the p. 19 characterisation (smallest
%         Q_p-subgroup whose extension to C contains the image of h).
%
% Announced and never established on the pages: the Question of p. 23 and the
% Question of p. 28 stay open; Gamma (p. 26) and S (pp. 27-28) are not defined;
% the fifth datum of p. 24 is abandoned.
\documentclass[11pt,a4paper]{article}
% Le préambule commun à toutes les transcriptions du fonds Grothendieck.
%
% Il tient en une idée : **l'incertitude se compose, elle ne se résout pas.**
% Une transcription qui lisse un mot douteux a détruit la seule chose pour
% laquelle on la faisait — un lecteur doit pouvoir savoir, à chaque endroit,
% ce qui était sur le feuillet et ce qui a été ajouté. Les six macros qui
% suivent sont donc l'appareil critique, et non de la décoration.
%
% Les mêmes commandes sont reconnues par `scripts/render.mjs`, qui produit la
% vue de lecture du site. Une macro ajoutée ici doit l'être aussi là-bas,
% faute de quoi le rendu échouera bruyamment — ce qui est voulu : un rendu
% silencieusement faux est pire qu'un rendu absent.

\ProvidesPackage{grothendieck}[2026/08/08 Transcriptions du fonds Grothendieck]

\usepackage{amsmath,amssymb,amsthm}
\usepackage{mathrsfs}
\usepackage{stmaryrd}
% Les diagrammes commutatifs sont partout dans ce fonds : c'est la langue
% même de la géométrie algébrique de Grothendieck, et une transcription qui
% les rendrait en prose ne transcrirait rien.
\usepackage{tikz-cd}
% Français par défaut — c'est la langue des feuillets — mais l'anglais est
% chargé aussi : la traduction et le résumé sont des documents anglais, et la
% typographie française (espaces avant les deux-points) y serait une faute.
% Ces documents basculent par \selectlanguage{english} en tête de corps.
\usepackage[english,french]{babel}
\usepackage{fontspec}
\usepackage{geometry}
\geometry{a4paper,margin=25mm,marginparwidth=28mm}
\usepackage{marginnote}
\usepackage{xcolor}
% ulem, not soul. `soul`'s \st and \ul are built on a character-by-character
% scan that XeLaTeX's font machinery breaks: the document fails with an
% undefined control sequence rather than degrading. ulem does the same two jobs
% and is Unicode-engine safe. `normalem` stops it hijacking \emph, which the
% transcriptions use for Grothendieck's own emphasis.
\usepackage[normalem]{ulem}
\usepackage{hyperref}
\hypersetup{colorlinks=true,linkcolor=black,urlcolor=black,pdfborder={0 0 0}}

% --- Métadonnées du lot -----------------------------------------------------
% Reprises telles quelles par le script de rendu, qui les affiche en tête de
% la vue de lecture. Elles fixent ce que le fichier prétend couvrir : sans
% elles, une transcription flotte sans point d'attache dans un fonds de seize
% mille feuillets.

\newcommand{\folder}[1]{\gdef\grothFolder{#1}}
\newcommand{\batch}[1]{\gdef\grothBatch{#1}}
\newcommand{\leaves}[2]{\gdef\grothFirst{#1}\gdef\grothLast{#2}}
\newcommand{\foldertitle}[1]{\gdef\grothFoldertitle{#1}}
\gdef\grothFolder{?}\gdef\grothBatch{?}\gdef\grothFirst{?}\gdef\grothLast{?}\gdef\grothFoldertitle{}

% --- Le résumé --------------------------------------------------------------
%
% Il ouvre la lecture modernisée, sous le titre « Résumé », et il est composé
% en petit. Ce n'est pas un ornement : c'est le seul passage du document qui
% ne soit pas des mathématiques, et un lecteur doit pouvoir le reconnaître
% avant de l'avoir lu — puis le sauter s'il connaît déjà le sujet, ou s'y
% arrêter s'il ne le connaît pas. Le filet qui le suit marque la reprise du
% corps.

\newenvironment{resume}
  {\small\setlength{\parskip}{0.45em}}
  {\par\medskip\hrule\medskip}

% --- Mots-clés ---------------------------------------------------------------
%
% En anglais, en clôture du résumé de la lecture modernisée : le vocabulaire
% moderne sous lequel chercher ce que les feuillets construisent. Le manifeste
% du site (`npm run manifest`) les extrait pour étiqueter la cote entière —
% cette ligne est la source unique des tags, il n'y a pas de fichier à côté.
\newcommand{\keywords}[1]{\par\smallskip\noindent
  {\footnotesize\textsc{Keywords} --- \textit{#1}}\par}

% --- L'appareil critique ----------------------------------------------------

% Le numéro de feuillet, en marge. C'est l'ancre qui permet de régler un
% désaccord sur une formule en regardant *un* feuillet plutôt qu'une chemise
% de six cents. Le site s'en sert aussi pour tourner les pages du fac-similé
% quand on fait défiler la transcription.
\newcommand{\leaf}[1]{%
  \marginnote{\footnotesize\color{gray!70}\textsf{#1}}%
  \label{leaf:#1}\ignorespaces}

% Illisible : on ne devine pas. Un mot inventé qui se lit comme les autres est
% le pire résultat possible.
\newcommand{\ill}{\textcolor{red!60!black}{[\,\ldots\,]}}

% Lecture douteuse : le mot est donné, et signalé comme incertain.
\newcommand{\uncertain}[1]{\uline{#1}}

% Ajout de l'éditeur — une lettre manquante, un mot sous-entendu.
\newcommand{\add}[1]{\textcolor{blue!50!black}{[#1]}}

% Biffé par Grothendieck lui-même. Ce qu'il a rayé fait partie du document :
% une rature dit souvent où la pensée a changé de direction.
\newcommand{\struck}[1]{\sout{#1}}

% Note du transcripteur, sur l'état matériel du feuillet.
\newcommand{\note}[1]{\par\smallskip{\small\color{gray!60!black}\textit{[#1]}}\par\smallskip}

% Note marginale de Grothendieck, distincte du corps du texte.
\newcommand{\marginal}[1]{\par\smallskip\noindent\hspace{1em}%
  {\small\color{gray!60!black}#1}\par\smallskip}

% --- Titre ------------------------------------------------------------------

\AtBeginDocument{%
  \begin{center}
    {\large\bfseries Fonds Alexandre Grothendieck --- cote n\textsuperscript{o}~\grothFolder}\\[2pt]
    {\itshape\grothFoldertitle}\\[4pt]
    {\small feuillets \grothFirst--\grothLast\ (lot \grothBatch)}\\[2pt]
    {\footnotesize\color{gray!60!black}Transcription établie d'après le fac-similé numérisé
     par l'Université de Montpellier.\\
     Les lectures incertaines sont soulignées, les illisibles notées [\,\ldots\,],
     les ajouts entre crochets.}
  \end{center}
  \vspace{1em}}

\endinput


\folder{161-5}
\pages{1}{28}
\dating{1967-1969 — le groupe « Dossiers rassemblés par Grothendieck sur différentes thématiques » (161-1 à 162-6) est daté [après 1961-vers 1977]}
\watermark{Édition de démonstration\\interprétation personnelle de l'œuvre}
\foldertitle{Barsotti-Tate : tirés à part annotés (1967-1969), notes manuscrites (s.d.) — lecture modernisée du dossier entier}

\begin{document}

\section*{Résumé}

\begin{resume}
La chemise porte un seul mot, « Barsotti-Tate », et contient deux choses de
nature différente : deux tirés à part imprimés de Jean-Pierre Serre, annotés à
l'encre rouge par une main qu'on ne sait pas attribuer, puis sept feuillets de
la main de Grothendieck. Les uns et les autres portent sur la même question :
comment comparer deux descriptions d'un même objet arithmétique, l'une par ses
symétries galoisiennes, l'autre par de l'algèbre linéaire.

Le point de départ est une situation que l'on connaît sur les nombres
complexes. La cohomologie d'une variété complexe porte deux structures : un
réseau, formé des cycles qu'on peut compter, et une décomposition de Hodge,
qui vient des formes différentielles. On passe de l'une à l'autre en intégrant
une forme le long d'un cycle ; les nombres obtenus sont les périodes. Sur un
corps $p$-adique, les deux structures existent encore. D'un côté, une
représentation du groupe de Galois, fabriquée avec les points de torsion. De
l'autre, un espace vectoriel muni d'un Frobenius et d'une filtration, qu'on lit
sur la réduction modulo $p$. Mais il n'y a plus d'intégrale pour passer de
l'une à l'autre. Les objets sur lesquels la question se pose d'abord sont les
groupes de Barsotti-Tate, le nom que Grothendieck donne aux groupes
$p$-divisibles : des tours de points de torsion, comme celles des racines
$p^n$-ièmes de l'unité. Tate a montré (1967) que le côté galoisien détermine
le groupe à isogénie près, et qu'il se décompose en deux morceaux après
extension des scalaires au complété de la clôture algébrique ; c'est la
décomposition de Hodge-Tate, l'ombre $p$-adique de celle de Hodge.

L'idée des feuillets est d'écrire chaque structure comme un \emph{torseur}.
Un torseur, c'est l'ensemble de toutes les façons d'identifier deux objets
isomorphes sans l'être canoniquement : l'ensemble des bases d'un espace
vectoriel en est le modèle, où aucune base n'est privilégiée mais où deux
bases quelconques diffèrent d'une unique matrice de passage. La
représentation galoisienne, la graduation de Hodge-Tate, la filtration, le
Frobenius deviennent autant de torseurs sous un même groupe algébrique, et les
feuillets en dressent la liste --- deux fois, la première abandonnée en
route. Puis ils posent la question de l'unicité : une fois donnée
l'action du groupe de Galois, y a-t-il, à conjugaison près, une seule façon de
lui adjoindre le Frobenius ? Elle devient une question de cohomologie
galoisienne non commutative. Telle qu'elle est posée, la réponse est non, et
la raison en est instructive : la liste ne relie le Frobenius à la
représentation galoisienne qu'à travers la graduation de Hodge-Tate, et cela
ne suffit pas.

Les deux imprimés de Serre sont le cadre de ce travail, et l'un d'eux le dit
lui-même : une note de bas de page imprimée y rapporte que l'idée de remplacer
les algèbres de Lie $p$-adiques par des groupes algébriques venait de
Grothendieck. Les annotations en rouge renvoient à des résultats de S. Sen et
à un livre de T. Springer, qui semblent postérieurs de quelques années aux
tirés à part ; c'est une inférence, non une date.

La question que posent ces feuillets est, sous une forme tannakienne, celle
que Grothendieck a posée publiquement en 1970 sous le nom de problème du
« foncteur mystérieux » ; les feuillets n'emploient pas l'expression, et rien
n'y dit s'ils sont d'avant ou d'après. C'est Jean-Marc Fontaine qui y a
répondu, à partir de la fin des années 1970, en construisant des anneaux de
périodes $p$-adiques qui jouent le rôle de l'intégrale manquante ; les
feuillets ne citent pas ces travaux. Les noms à
chercher ensuite sont la théorie de Hodge $p$-adique, les représentations
cristallines, les $\varphi$-modules filtrés. Le bas de la dernière page, enfin,
est un calcul séparé sur un tore et son modèle de Néron.

\keywords{p-adic Hodge theory, Hodge–Tate representation, Barsotti–Tate group, Tannakian category, filtered fibre functor, F-isocrystal, mysterious functor, non-abelian Galois cohomology, Sen theory, p-adic Mumford–Tate group, Néron model}
\end{resume}

\pagerange{1}{28}
\section*{Le fil du dossier, et les conventions}

\emph{Les campagnes.} Le dossier réunit des pièces de dates et de natures
différentes.

\begin{enumerate}
  \item \emph{Page 1}, la chemise : le seul mot « Barsotti-Tate », au crayon,
  sur l'onglet.
  \item \emph{Pages 2 à 16}, imprimé : J.-P. Serre, « Sur les groupes de
  Galois attachés aux groupes $p$-divisibles », exposé de la conférence de
  Driebergen (1966), paru chez Springer en 1967. Annotations aux pages 9, 15
  et 16.
  \item \emph{Pages 17 à 21}, imprimé : le résumé du cours de Serre au Collège
  de France pour 1967-1968, « Algèbre et géométrie », extrait de l'Annuaire.
  Annotations aux pages 19 et 20 ; la page 21 ne porte rien de manuscrit.
  \item \emph{Pages 22 et 23}, encre sur papier quadrillé : le cadre, les
  données cristallines, une première Question.
  \item \emph{Pages 24 à 26}, crayon sur papier réglé : la liste des données,
  commencée deux fois --- la page 24 s'arrête à la cinquième donnée, la page
  25 recommence à la première, et la page 26 poursuit la page 25.
  \item \emph{Pages 27 et 28}, crayon sur papier quadrillé : la cohomologie
  non abélienne de l'extension, et une seconde Question.
  \item \emph{Bas de la page 28}, encre : un calcul indépendant sur les tores
  d'un corps $p$-adique.
\end{enumerate}

\emph{Le fil.} Les tirés à part donnent le cadre : l'enveloppe algébrique de
l'image de Galois, la décomposition de Hodge-Tate, l'homomorphisme $h$ de
$\mathbb{G}_m$ qui la code. Les feuillets manuscrits reprennent ce cadre sous
forme tannakienne --- un groupe pro-algébrique $G$, une représentation
$\rho$ de $\pi$ dans $G$, un cocaractère $j$ --- et y ajoutent ce que Serre
n'a pas : les structures qui viennent de la réduction modulo $p$, la
filtration de de Rham et le Frobenius. La Question de la page 23 demande si
ces structures, que la page tient pour acquises sur les groupes de
Barsotti-Tate, se prolongent à
toute la catégorie tensorielle qu'ils engendrent. Les pages 24 à 26 cherchent
à dire ce que serait une telle structure pour un groupe $G$ quelconque ; les
pages 27 et 28 demandent si elle est unique. Il n'y a pas de démonstration
dans le dossier : il y a un cadre, une liste, et deux questions.

\emph{Notations.} Elles valent pour tout le document, et diffèrent par
endroits de celles des pages ; chaque écart est signalé en note au point où il
se produit.

\begin{itemize}
  \item $K$ est un corps de caractéristique $0$, complet pour une valuation
  discrète, de corps résiduel $k$ \emph{parfait} de caractéristique $p$ ;
  $\mathcal{O}_K$ son anneau d'entiers, $W = W(k)$ l'anneau des vecteurs de
  Witt, $K_0 = W[1/p] \subset K$, et $\sigma$ le Frobenius de $K_0$.
  \item $\bar{K}$ une clôture algébrique de $K$, $\Omega$ son complété,
  $\pi = \mathrm{Gal}(\bar{K}/K)$ et $\pi' = \mathrm{Gal}(\bar{K}/K_0)$, qui
  contient $\pi$ avec l'indice $[K : K_0]$.
  \item $\Omega(i) = \Omega \otimes_{\mathbb{Z}_p} \mathbb{Z}_p(1)^{\otimes i}$,
  où $\mathbb{Z}_p(1)$ est le module de Tate des racines de l'unité de
  $\bar{K}$ ; $\pi$ y agit diagonalement.
  \item Pour une représentation $M$ : $\omega(M) = M \otimes_{\mathbb{Q}_p}
  \Omega$ avec $\pi$ agissant sur $\Omega$ seul, et $\omega^{!}(M)$ le même
  espace avec l'action diagonale.
\end{itemize}

\emph{Deux lettres qui changent de rôle.} La page 22 pose
$L = \mathrm{Frac}\, W$, contenu dans $K$. À partir du point c) de la page 23
et jusqu'à la fin, les lettres $K$ et $L$ sont échangées : $L$ y désigne le
corps de base et $K$ le sous-corps absolument non ramifié.\footnote{Les
indices sont nombreux et concordent. Page 23 :
$T_{\mathrm{DR}}(M) = T_{\mathrm{cr}}(M) \otimes_K L$, alors que
$T_{\mathrm{cr}}(M)$ vient d'être défini comme un $L$-vectoriel. Page 24 :
$Q_L$ « descendant $Q_{\Omega}$ », là où la page 22 descendait le torseur de
Hodge-Tate « sur $K$ ». Page 25 : $\pi = \mathrm{Aut}_L(\Omega)$, et l'Axiome descend le
torseur de Hodge-Tate « à $L$ ».
Page 26 : le foncteur de de Rham est un « $L$-foncteur », la descente va « de
$L : K$ », et la donnée 6) est un isomorphisme de $P_K$ avec un tordu
$P_K^{(\sharp)}$ qui ne peut guère être qu'un tordu par le Frobenius, ce qui
n'a de sens que si $K$ porte un Frobenius, donc est non ramifié. Pages 27 et 28 : $G'_S(S)^{\pi} = G'_L(L)$, ce qui exige que $L$ soit
le corps des invariants de $\pi$. Lu avec la convention de la page 22, chacun
de ces passages est incorrect ou dépourvu de sens ; lu avec l'échange, tous
sont justes. La
transcription relève l'échange tel que les pages l'écrivent, en une seule note
au point c) de la page 23, et garde partout ses lettres.} On écrit ici $K$ et $K_0$ partout.

\pagerange{1}{21}
\section*{I. Les deux tirés à part de Serre et leurs annotations (pages 1 à 21)}

Les pages 2 à 21 sont imprimées, et le texte imprimé n'est pas de lui ; il
n'est pas transcrit, et on ne le résume ici que dans la mesure où il faut
situer une annotation. Ce qui appartient au dossier, ce sont cinq annotations
à l'encre rouge, toutes de la même main. \emph{Cette main n'a pas pu être
identifiée avec certitude comme celle de Grothendieck} : rien dans l'écriture
n'est assez caractéristique pour trancher, et ce qui suit parle donc
d'« annotations », jamais de « ses » annotations.\footnote{Il en va de même du
mot « Barsotti-Tate » écrit au crayon, en long, sur l'onglet de la chemise
(page 1) : la main n'en est pas sûrement la sienne. Une sixième intervention,
un signe $<$ repassé à l'encre bleue sur une retouche blanche, page 15, dans
la démonstration du lemme 3 de l'article, n'est attribuée à personne ; ce
qu'elle recouvre ne se lit pas.}

Il faut dire pourquoi ces imprimés sont là. Le premier porte, à sa page 3, une
note de bas de page imprimée où Serre écrit que l'idée de remplacer les
algèbres de Lie $p$-adiques par les groupes algébriques correspondants lui
avait été suggérée par Grothendieck, « en liaison avec sa théorie des
motifs ». C'est le texte de Serre, non une annotation, et il n'est pas
transcrit ; mais c'est le geste même que font les feuillets manuscrits des
pages 22 et suivantes, où le groupe de Galois est remplacé par son enveloppe
algébrique $G$.

\pagerange{2}{16}
\subsection*{« Sur les groupes de Galois attachés aux groupes $p$-divisibles » (pages 2 à 16)}

L'article, paru dans les actes de la conférence de Driebergen sur les corps
locaux (Springer, 1967) --- le volume même où Tate publie son article
« $p$-divisible groups », sur lequel reposent les pages 22 et 23 ---, étudie
l'image du groupe de Galois dans le module de Tate d'un groupe $p$-divisible.
Trois de ses pages sont annotées.

\emph{Page 9.} En regard de la remarque 3 qui clôt le § 3, où l'imprimé
établit une inclusion $M \subset I_{\mathrm{alg}}$ entre le groupe de
Mumford-Tate $M$ du module galoisien $V$ et l'enveloppe algébrique
$I_{\mathrm{alg}}$ de l'image de l'inertie, la marge porte : « En fait, on a
$M = I_{\mathrm{alg}}$ (Sen) ». L'égalité ainsi affirmée demande une
précision : elle vaut pour la composante neutre.\footnote{La note de la
transcription à la page 9 dit seulement que $M$ est « le groupe de
Mumford-Tate du module galoisien $V$ ». La caractérisation qui suit est celle
que le cours de la page 19 donne de $H_V^{\circ}$, l'analogue $p$-adique du
groupe de Mumford-Tate, et que rapporte la note de la transcription à cette
page. Si $M$ est ainsi le plus petit sous-groupe algébrique sur
$\mathbb{Q}_p$ dont l'extension à $C$, le complété de $\bar{K}$, contient
l'image de $h$, il est connexe, alors
que $I_{\mathrm{alg}}$ ne l'est pas toujours : un caractère ramifié d'ordre
fini est de Hodge-Tate de poids $0$, son $M$ est trivial et son
$I_{\mathrm{alg}}$ est un groupe fini non trivial. Ce que le théorème de Sen
cité plus bas permet d'établir est $M = (I_{\mathrm{alg}})^{\circ}$ :
l'algèbre de Lie de l'image de l'inertie est algébrique, c'est celle de $M$,
et un sous-groupe ouvert des points $p$-adiques d'un groupe connexe y est
Zariski-dense. On n'a pas vu l'imprimé et l'on ne peut pas dire si Serre y
définit $M$ autrement.}

\emph{Page 15.} Contre la remarque 1 qui suit le théorème 5 --- sans
l'hypothèse $\mathrm{End}(F) = \mathbb{Z}_p$ sur le groupe formel $F$ de
dimension 1, l'algèbre de Lie de l'image de Galois est encore déterminée,
par une démonstration que l'imprimé dit « sensiblement plus compliquée » ---
la marge porte : « voir Sen (Annals) ».\footnote{Le dernier mot est écrit vite
et abrégé ; la lecture « Annals » est celle de la page 19, où il se lit en
clair. L'annotation ne dit pas si Sen démontre la remarque ou en simplifie la
preuve.}

\emph{Page 16.} Un crochet embrasse la question qui termine l'article. Le
groupe d'inertie $I$ y porte deux filtrations, l'une par la congruence de
l'action sur le module de Tate $T$, $v(g) \geqslant n \iff (g-1)(T) \subset
p^n T$, l'autre par les groupes de ramification en numérotation supérieure,
$w(g) \geqslant n \iff g \in I^n$ ; l'imprimé demande si $w = e_K\, v +
O(1)$, $e_K$ désignant l'indice de ramification absolu, et pose la même
question pour tout groupe de Galois qui est un groupe de Lie $p$-adique. La
marge porte : « démontré par S. Sen (Inventions) ». C'est exact si, comme il
est vraisemblable, l'annotation vise l'article de Sen sur la ramification des
extensions de Lie $p$-adiques (Inventiones Math. 17, 1972), dont le théorème
compare précisément la filtration de ramification supérieure et la filtration
de Lie, à une constante près, pour toute extension galoisienne totalement
ramifiée de groupe de Lie $p$-adique.\footnote{L'identification des
références est la nôtre, et celle de la transcription : les annotations
écrivent seulement « Annals » et « Inventions ».}

\pagerange{17}{21}
\subsection*{« Algèbre et géométrie », résumé du cours de 1967-1968 (pages 17 à 21)}

L'extrait de l'Annuaire du Collège de France (68e année, 1968-1969) résume un
cours sur les représentations $p$-adiques de Hodge-Tate. Une représentation
$V$ du groupe de Galois y a pour image $G_V$ ; $H_V$ est l'enveloppe
algébrique de $G_V$, et un homomorphisme $h$ de $\mathbb{G}_m$ dans
$\mathrm{GL}_V$, défini sur le complété $C$ de la clôture algébrique, code la
décomposition de Hodge-Tate. Deux pages sont annotées ; la dernière page
imprimée, la page 21, ne porte rien de manuscrit.

\emph{Page 19.} L'imprimé caractérise, lorsque $G_V$ est résoluble, la
composante neutre $H_V^{\circ}$ comme le plus petit sous-groupe algébrique de
$\mathrm{GL}_V$ sur $\mathbb{Q}_p$ dont l'extension à $C$ contient l'image de
$h$ --- l'analogue $p$-adique du groupe de Mumford-Tate ---, conjecture que la
caractérisation vaut sans l'hypothèse de résolubilité, et observe qu'il
suffirait de savoir qu'une représentation de Hodge-Tate dont tous les poids
sont nuls a une image finie. La marge porte, en face de cet alinéa : « ceci a
été démontré par S. Sen (Annals) ». C'est exact, pour l'image de l'inertie :
l'article de Sen sur les algèbres de Lie des groupes de Galois attachés aux
modules de Hodge-Tate (Annals of Math. 97, 1973) établit que l'algèbre de Lie
de l'image de l'inertie est la plus petite algèbre de Lie algébrique sur
$\mathbb{Q}_p$ dont l'extension à $C$ contient celle de l'image de $h$, et du
même coup qu'une représentation de Hodge-Tate de poids tous nuls a une inertie
d'image finie. Lorsque le corps résiduel est algébriquement clos, l'inertie
est tout le groupe, et c'est la conjecture du cours.\footnote{Pour un corps
résiduel fini, l'énoncé ne peut porter que sur l'inertie : un caractère non
ramifié d'ordre infini est de Hodge-Tate de poids nul et d'image infinie. Le
cadre exact de l'imprimé n'est pas visible dans la transcription, qui ne le
reproduit pas. Même réserve sur l'identification de la référence que
ci-dessus. Le rapprochement avec les feuillets manuscrits, où le groupe $G_M$
de la page 22 est le $H_V$ de Serre, est le nôtre.}

\emph{Page 20.} Le point 5) du résumé traite le cas où $V$ n'a que les poids
$0$ et $1$, qui est celui du module d'un groupe $p$-divisible, le nombre $n_1$
de poids égaux à $1$ étant la dimension du groupe. Pour $n_1 > 1$, dit la
dernière phrase, le cours s'est achevé sur un essai de classification fondé
sur ce que $H_V^{\circ}$ contient un tore de dimension 1 n'ayant que deux
poids. La marge porte : « voir à ce sujet le livre de T. Springer sur les
alg. de Jordan récemment paru chez Springer-Verlag ».\footnote{Le livre visé
est vraisemblablement T. A. Springer, \emph{Jordan algebras and algebraic
groups}, Springer, 1973. Le lien qu'on devine est celui-ci : un cocaractère
qui n'a que deux poids sur $V$ --- on dit aujourd'hui \emph{minuscule} ---
induit sur l'algèbre de Lie une graduation en degrés $-1$, $0$, $1$, et de
telles graduations sont liées aux structures de Jordan. C'est notre lecture de
l'annotation ; nous n'affirmons pas que le livre traite ce cas. La
classification des groupes de monodromie munis d'un tel cocaractère a été
faite depuis, notamment par R. Pink (1998).}

\emph{Ce que les annotations disent de leur date.} Si les références sont
bien celles qu'on vient de dire, les annotations sont postérieures à 1972, et
celle de la page 20 --- « récemment paru » --- plutôt de 1973 ou peu après :
plus tard, donc, que les années 1967-1969 que l'inventaire donne aux
tirés à part. C'est une inférence, non une donnée des pages, et la datation
de Montpellier est reproduite telle quelle en tête de ce document. Les
feuillets manuscrits, eux, ne portent aucune date.\footnote{Un indice, faible,
va dans le même sens pour eux : le nom de « groupe de Barsotti-Tate » est
celui que Grothendieck donne aux groupes $p$-divisibles, et sous lequel il en
traite publiquement en 1970, au congrès de Nice et dans son cours de
Montréal. Cela n'exclut pas qu'il l'ait employé plus tôt.}

\pagerange{22}{23}
\section*{II. Le cadre : représentations de Hodge-Tate et groupes de Barsotti-Tate (pages 22 et 23)}

Les deux premiers feuillets, à l'encre, posent le cadre avec soin ; c'est la
seule partie du dossier qui ait l'allure d'une mise au net.

\subsection*{Les corps}

On se donne $K$ comme dans les conventions : caractéristique $0$, complet
pour une valuation discrète, corps résiduel $k$ parfait de caractéristique
$p$ --- la page souligne l'hypothèse d'un point d'exclamation.\footnote{La
page dit « $K$ corps local d'inégales caractéristiques » et note $V$ l'anneau
des entiers, $W$ le « sous-$p$-anneau de corps résiduel $k$ », $L$ son corps
des fractions. Pour $k$ parfait, ce sous-anneau est l'anneau des vecteurs de
Witt $W(k)$, et $L$ est le corps que l'on note aujourd'hui $K_0$ ; on garde
$W$, et l'on écrit $\mathcal{O}_K$ et $K_0$, la lettre $V$ servant plus haut
aux représentations et la lettre $L$ changeant de sens dans la suite des
feuillets. La première ligne de la page est biffée et illisible.} Les
inclusions, que la page écrit en marge, sont
\[
\begin{array}{ccccccc}
K_0 & \subset & K & \subset & \bar{K} & \subset & \Omega \\
\cup & & \cup & & \cup & & \cup \\
W & \subset & \mathcal{O}_K & \subset & \mathcal{O}_{\bar{K}} & \subset & \mathcal{O}_{\Omega}
\end{array}
\]
avec $\pi = \mathrm{Gal}(\bar{K}/K)$ et $\pi' = \mathrm{Gal}(\bar{K}/K_0)$. Le
groupe $\pi'$ est introduit, suivi de points de suspension, et ne sert pas
sur ces deux pages ; on le retrouvera page 26. Le fait de base, que les pages
utilisent sans le rappeler, est que $\Omega^{\pi} = K$ : c'est le théorème de
Tate (1967), dit aujourd'hui d'Ax-Sen-Tate.

\subsection*{La catégorie des représentations de Hodge-Tate}

Un $\mathbb{Q}_p$-faisceau lisse sur $\mathrm{Spec}\, K$ est la même chose
qu'une représentation continue $M$ de $\pi$ sur un $\mathbb{Q}_p$-espace de
dimension finie. Elle est dite \emph{de Hodge-Tate} si $\omega^{!}(M) = M
\otimes_{\mathbb{Q}_p} \Omega$, muni de l'action diagonale, est isomorphe
comme $\Omega$-espace semi-linéaire à une somme $\bigoplus_i
\Omega(i)^{h_i}$ ; les entiers $i$ tels que $h_i \neq 0$ sont ses
\emph{poids}, et la convention est celle où $\mathbb{Q}_p(1)$ est de poids
$1$.\footnote{La page parle des degrés de « $\mathrm{gr}(M_{\Omega})$ » : son
$M_{\Omega}$ est notre $\omega^{!}(M)$, l'action diagonale étant la seule pour
laquelle cette graduation a un sens. La page 25 écrit $M^{!}_{\Omega}$ pour
le même objet et réserve $M_{\Omega}$ à l'action sur $\Omega$ seul ; on suit
ici la page 25.}

On note $\mathcal{M}_K$ la catégorie des représentations de Hodge-Tate. Elle
est stable par sommes, produits tensoriels, duaux et sous-quotients, et le
foncteur qui oublie l'action de $\pi$ en fait une catégorie tannakienne
neutre sur $\mathbb{Q}_p$ --- la page dit « $\otimes$-catégorie sur
$\mathbb{Q}_p$ ». Pour $n \geqslant 0$, on note $\mathcal{M}^{(n)}_K$ la
sous-catégorie pleine des représentations dont les poids sont compris entre
$0$ et $n$ ; leur réunion $\mathcal{M}^{+}_K$ est la sous-catégorie des
objets \emph{effectifs}. La sous-catégorie $\mathcal{M}^{(1)}_K$ est celle des
représentations de poids $0$ et $1$.\footnote{La page définit une seule
catégorie, dont l'exposant se lit « $(n)$ », comme formée des $M$ « avec
$\mathrm{gr}\, M_{\Omega}$ à degrés $0$ et $1$ », et la ligne suivante
emploie $\mathcal{M}^{(1)}_K$ ; l'accolade ouverte n'est pas refermée. La
définition de $\mathcal{M}^{(n)}_K$ par les poids de $0$ à $n$, qui rend
correcte l'écriture $\mathcal{M}^{+}_K = \varinjlim \mathcal{M}^{(n)}_K$ de la
ligne précédente, est la nôtre.}

\subsection*{Le plongement de Tate}

Soit $\widetilde{\mathrm{BT}}$ la catégorie des groupes de Barsotti-Tate sur
$\mathcal{O}_K$ à isogénie près : mêmes objets que les groupes
$p$-divisibles sur $\mathcal{O}_K$, morphismes
$\mathrm{Hom}(G, H) \otimes \mathbb{Q}$.\footnote{« Groupe de Barsotti-Tate »
est le nom que Grothendieck donne aux groupes $p$-divisibles de Tate. Sur la
page, l'indice $\mathcal{O}_K$ (son $V$) est écrit sur un $K$ biffé.} À un tel
groupe on associe $V_p(G) = T_p(G) \otimes_{\mathbb{Z}_p} \mathbb{Q}_p$,
représentation de $\pi$. Deux théorèmes de Tate, publiés dans le même volume
de Driebergen que le premier tiré à part, donnent alors :

\begin{itemize}
  \item le foncteur $G \mapsto V_p(G)$ de $\widetilde{\mathrm{BT}}$ dans les
  représentations de $\pi$ est pleinement fidèle ;
  \item $V_p(G)$ est de Hodge-Tate, de poids $1$ avec multiplicité
  $\dim G$ et de poids $0$ avec multiplicité $\mathrm{ht}\, G - \dim G$.
\end{itemize}

On obtient ainsi le plongement pleinement fidèle de la page,
$\widetilde{\mathrm{BT}} \hookrightarrow \mathcal{M}^{(1)}_K$, et l'on
identifie $\widetilde{\mathrm{BT}}$ à une sous-catégorie pleine de
$\mathcal{M}^{(1)}_K$, comme le dit la marge. Elle est \emph{strictement}
plus petite : la représentation de Tate d'une courbe de Tate sur $K$ est de
poids $0$ et $1$ et ne vient d'aucun groupe de Barsotti-Tate sur
$\mathcal{O}_K$.\footnote{On le sait aujourd'hui parce qu'elle n'est pas
cristalline, alors que les $V_p(G)$ le sont ; ce critère vient des travaux de
Fontaine, que les pages ne citent pas. L'exemple est le nôtre.}

On note $\mathcal{M}^{\mathrm{BT}}_K$ la sous-catégorie tannakienne de
$\mathcal{M}_K$ engendrée par l'image de $\widetilde{\mathrm{BT}}$ : la plus
petite sous-catégorie pleine qui la contient et qui est stable par sommes,
produits tensoriels, duaux et sous-quotients.\footnote{La page écrit
« catégorie tannakienne engendrée par $\mathcal{M}^{(1)}_K$ ». Avec cette
définition, la Question de la page 23, qui demande de \emph{prolonger} à
$\mathcal{M}^{\mathrm{BT}}$ des données définies sur les groupes de
Barsotti-Tate, porterait aussi sur des objets de $\mathcal{M}^{(1)}_K$ qui ne
viennent d'aucun groupe, comme celui de la courbe de Tate ; on suit l'intention
de la Question. L'en-tête de la transcription dit d'ailleurs, lui aussi, « la
sous-catégorie tannakienne qu'ils engendrent », en parlant des groupes.}

\subsection*{Le groupe $G$}

Le foncteur qui à $M$ associe son $\mathbb{Q}_p$-espace sous-jacent est noté
$T_p$ par la page --- pour $M = V_p(G)$, c'est le module de Tate rationnel ---
et appelé « foncteur fibre sur $\Omega$ » : c'est la fibre du faisceau $M$ au
point géométrique défini par $\Omega$. Son groupe d'automorphismes tensoriels
\[
G = \mathrm{Aut}^{\otimes}(T_p)
\]
est un groupe pro-algébrique sur $\mathbb{Q}_p$, et $\mathcal{M}_K$ est
équivalente à la catégorie de ses représentations de dimension finie. L'action
de $\pi$ fournit un homomorphisme continu $\rho : \pi \to G(\mathbb{Q}_p)$ ;
son image est Zariski-dense, et c'est en ce sens que $G$ est l'\emph{enveloppe
algébrique} de $\pi$.\footnote{La page ajoute « NB $G/G^{\circ}$ non », suivi
d'un mot illisible ; on ne le restitue pas. Le second $G$ de la ligne portait
un exposant, biffé ; la page 25 dit que $G$ dépend en fait du choix de
$\Omega$, ce que l'exposant notait sans doute.} Pour chaque $M$, l'image
\[
G_M = \mathrm{Im}\bigl(G \to \mathrm{GL}(T_p(M))\bigr)
\]
est l'adhérence de Zariski de l'image de $\pi$ dans $\mathrm{GL}(T_p(M))$ :
c'est le groupe que le cours de Serre appelle $H_V$. Remplacer l'image de
Galois, groupe de Lie $p$-adique, par son enveloppe algébrique, c'est l'idée
que la note imprimée de Serre attribue à Grothendieck ; passer des groupes
$H_V$, un par représentation, au seul groupe $G$ qui les domine tous, c'est
passer au langage tannakien.\footnote{« Catégorie
tannakienne » est un nom de Grothendieck lui-même, et les pages l'emploient ;
la théorie en a été rédigée par son élève N. Saavedra Rivano,
\emph{Catégories tannakiennes} (1972).}

\subsection*{La structure de Hodge-Tate, écrite en torseurs}

Le cœur des deux pages est une réécriture de la décomposition de Hodge-Tate
en termes de torseurs, en trois étages que la page dispose en tableau : une
structure « sur $(\Omega, \pi)$ », la même tordue, puis sa descente sur
$K$.\footnote{Dans le tableau de la page, chaque ligne porte une étiquette, un
morphisme de $\mathbb{G}_m$ dessiné en diagonale montante, un torseur et un
groupe. L'étiquette de la deuxième ligne se lit « id. », sous réserve. Ce qui
suit explicite le tableau à l'aide de ce que la page 25 dit des mêmes objets ;
la présentation par les deux foncteurs fibres est celle de la page 25.}

On travaille dans la catégorie des $\Omega$-espaces de dimension finie munis
d'une action semi-linéaire continue de $\pi$. Les deux foncteurs $\omega$ et
$\omega^{!}$ des conventions y envoient les représentations de $G$, et ce sont
deux foncteurs fibres. Le torseur
\[
\mathfrak{P} = \mathrm{Isom}^{\otimes}(\omega, \omega^{!})
\]
est un torseur à droite sous $\mathrm{Aut}^{\otimes}(\omega) = G_{\Omega}$, où
$\pi$ agit sur les coefficients. Il a une section évidente $e$, l'identité des
espaces sous-jacents, et cette section n'est pas invariante sous $\pi$ : on a
$\gamma(e) = e \cdot \rho(\gamma)$, de sorte que $\mathfrak{P}$ est le torseur
défini par le $1$-cocycle $\rho$.\footnote{« non » est une lecture incertaine,
que la formule justifie : la section $e$ n'est invariante que si $\rho$ est
trivial.} Son groupe d'automorphismes est $\mathrm{Aut}^{\otimes}(\omega^{!})
= G^{!}_{\Omega}$, c'est-à-dire $G_{\Omega}$ sur lequel $\gamma \in \pi$ agit
par son action sur $\Omega$ composée avec $\mathrm{int}(\rho(\gamma))$.

La décomposition de Hodge-Tate est une graduation du foncteur $\omega^{!}$ :
$\omega^{!}(M) = \bigoplus_i \omega^{!}(M)_i$, avec $\omega^{!}(M)_i \simeq
\Omega(i)^{h_i}$. Elle est unique et fonctorielle, parce que
$\mathrm{Hom}_{\pi}(\Omega(i), \Omega(i')) = 0$ pour $i \neq i'$ --- c'est
encore le théorème de Tate ---, et compatible au produit tensoriel. Une
graduation tensorielle d'un foncteur fibre, c'est un cocaractère de son
groupe ; on obtient ainsi
\[
j^{!} : \mathbb{G}_{m,\Omega} \longrightarrow G^{!}_{\Omega},
\]
où $t$ agit sur $\omega^{!}(M)_i$ par $t^i$. Il commute à $\pi$, parce que la
graduation est stable sous l'action diagonale. C'est l'homomorphisme $h$ du
cours de Serre, vu comme cocaractère d'un seul groupe.

Reste à rendre chaque morceau trivial. On pose
\[
\mathrm{Hdg}_{\Omega}(M) = \bigoplus_i \omega^{!}(M)_i \otimes_{\Omega}
\Omega(-i),
\]
qui est isomorphe, avec son action de $\pi$, à une somme de copies de
$\Omega$. Le torseur correspondant,
$Q_{\Omega} = \mathrm{Isom}^{\otimes}(\omega, \mathrm{Hdg}_{\Omega})$, se
déduit de $\mathfrak{P}$ en le tordant, à travers $j^{!}$, par le
$\mathbb{G}_m$-torseur $\mathcal{T}$ des générateurs de $\Omega(1)$ --- la
page écrit $Q_{\Omega} = \mathcal{T} \wedge^{\mathbb{G}_{m,\Omega}}
\mathfrak{P}$, le produit contracté qu'on note aujourd'hui
$\times^{\mathbb{G}_m}$. Son groupe d'automorphismes $G^{\mathrm{Hdg}}_{\Omega}$
est la forme tordue correspondante de $G^{!}_{\Omega}$, et $j^{!}$ y induit un
cocaractère $j_{\Omega}$.\footnote{La page note $\mathcal{T}$ ce torseur à la
page 22 et à la page 25, et $\mathcal{T}^{*}$ à la page 24 ; selon qu'on fait
agir $\mathbb{G}_m$ par $t^i$ ou $t^{-i}$ sur le poids $i$, il faut tordre par
$\mathcal{T}$ ou par son inverse. Le choix est fixé ici par la formule de la
page 25, $\mathrm{gr}^i \otimes S_{\Omega}(-i)$, qui est celle qui rend les
morceaux triviaux ; la page écrit $S_{\Omega}(1)$ pour notre $\Omega(1)$ et le
définit comme $T_p(\Omega^{*}) \otimes_{\mathbb{Z}_p} \Omega$. La marge de la
page 22, qui définit $\mathcal{T}$, est en partie illisible. Enfin la page 22
écrit $G^{\mathrm{Hdg}}$ et les pages 24 à 26 $G''$ pour le même groupe ; on
garde le premier nom.}

Comme chaque $\mathrm{Hdg}_{\Omega}(M)$ est trivial, il descend, et de façon
unique, sur $\Omega^{\pi} = K$ : le foncteur $D \mapsto D \otimes_K \Omega$ des
$K$-espaces vectoriels vers les $\Omega$-espaces à action semi-linéaire de
$\pi$ est pleinement fidèle, d'image les objets triviaux. On obtient le
foncteur fibre gradué sur $K$
\[
T_{\mathrm{Hdg}}(M) = \mathrm{Hdg}_{\Omega}(M)^{\pi} = \bigoplus_i \bigl(M
\otimes_{\mathbb{Q}_p} \Omega(-i)\bigr)^{\pi},
\]
et, « par descente sur $K$ », le torseur $Q_K =
\mathrm{Isom}^{\otimes}(T_p \otimes K, T_{\mathrm{Hdg}})$ sous $G_K$, le groupe
$G^{\mathrm{Hdg}}_K = \mathrm{Aut}^{\otimes}(T_{\mathrm{Hdg}})$, forme
intérieure de $G_K$, et le cocaractère de graduation $j_K :
\mathbb{G}_{m,K} \to G^{\mathrm{Hdg}}_K$. La dernière formule du tableau,
\[
T_{\mathrm{Hdg}}(M) = Q_K \wedge^{G} (T_p(M) \otimes K),
\]
dit simplement que le foncteur de Hodge-Tate est le foncteur fibre associé au
torseur $Q_K$.\footnote{C'est le foncteur que Fontaine notera
$D_{\mathrm{HT}}(M) = (M \otimes B_{\mathrm{HT}})^{\pi}$, avec
$B_{\mathrm{HT}} = \bigoplus_{i \in \mathbb{Z}} \Omega(i)$ ; l'anneau et la
notation sont de la fin des années 1970 ; les pages ne les emploient pas.
L'exposant de
$T_{\mathrm{Hdg}}$ est biffé sur la page.}

\subsection*{Les données cristallines, et la Question}

Pour un groupe de Barsotti-Tate, la page énumère alors quatre données, qu'elle
tient pour acquises, « fonctorielles en $M$ » :\footnote{La page commence par
« Pour $M \in \mathrm{Ob}\,\widetilde{\mathrm{BT}}$, on a
$T_{\mathrm{Hdg}}(M)$ », biffé, suivi d'un mot illisible. Une première
rédaction de a) disait « un $L$-vectoriel filtré $T_{\mathrm{cr}}(M)$ et un
isomorphisme $\mathrm{gr}\, T_{\mathrm{DR}}(M)_K \simeq T_{\mathrm{Hdg}}(M)$ » ;
le mot « filtré » et l'isomorphisme sont biffés, et reportés en c) et d). En
b), le mot « donnée » est biffé au profit de « structure ».}

\begin{itemize}
  \item a) un $K_0$-espace vectoriel $T_{\mathrm{cr}}(M)$ ;
  \item b) une structure de F-isocristal sur $T_{\mathrm{cr}}(M)$ : une
  bijection $\sigma$-semi-linéaire $F$ ;
  \item c) une filtration de $T_{\mathrm{DR}}(M) = T_{\mathrm{cr}}(M)
  \otimes_{K_0} K$ ;
  \item d) un isomorphisme gradué $\mathrm{gr}\, T_{\mathrm{DR}}(M) \simeq
  T_{\mathrm{Hdg}}(M)$.
\end{itemize}

Le point c) est celui où les lettres commencent à s'échanger : la page écrit
$T_{\mathrm{cr}}(M) \otimes_K L$.\footnote{Voir la note sur l'échange de $K$
et $L$ dans « Le fil du dossier ». Seul $\otimes_{K_0} K$ a un sens ici :
$T_{\mathrm{cr}}(M)$ est un $K_0$-espace, et la filtration de Hodge vit sur
$K$.} Ce que la page désigne ainsi est, pour $M = V_p(G)$, le module de
Dieudonné de la fibre spéciale de $G$ tensorisé par $\mathbb{Q}$, avec son
Frobenius, et la filtration de Hodge de ce module étendu à $K$ ; d) est alors,
à dualité près, le théorème de Hodge-Tate de Tate, dont les deux morceaux
s'expriment par l'algèbre de Lie de $G$ et celle de son dual.\footnote{Sur la
variance : le module de Dieudonné usuel est contravariant en $G$, alors que
$V_p$ est covariant. La page dit seulement « fonctorielles en $M$ ». Pour que
$T_{\mathrm{cr}}$ puisse être un $\otimes$-foncteur sur
$\mathcal{M}^{\mathrm{BT}}$ au sens de la Question, on le prend covariant en
$M$, c'est-à-dire qu'on compose la théorie contravariante avec la dualité ;
c'est un choix, et il est le nôtre. Sur les dates de ces notions : la
théorie de Dieudonné sur un corps parfait est classique, le cristal de
Dieudonné d'un groupe de Barsotti-Tate et sa filtration de Hodge sont le
programme de Grothendieck lui-même autour de 1970, développé par W. Messing
(1972), et d) combine ce programme avec le théorème de Tate. Que les $V_p(G)$
et ces données se correspondent exactement est aujourd'hui un théorème, établi
par Fontaine dans des cas particuliers, puis en général par Faltings et
Breuil.}

\emph{Question} (page 23). Peut-on prolonger ces données de
$\widetilde{\mathrm{BT}}$ à $\mathcal{M}^{\mathrm{BT}}_K$, de sorte que
$T_{\mathrm{cr}}$ soit un $\otimes$-foncteur, que la structure b) soit
compatible au produit tensoriel et aux $\mathrm{Hom}$ internes, et de même c)
et d) ? Le prolongement de $T_{\mathrm{cr}}$ est-il unique à isomorphisme
unique près ?\footnote{Un début de phrase biffé, « A-t-on », précède la
Question.}

C'est, sous une forme tannakienne, la question que Grothendieck a posée
publiquement en 1970 comme celle du « foncteur mystérieux » : retrouver la
structure cristalline à partir de la représentation galoisienne. La réponse à
la première partie est oui, et elle est venue d'ailleurs : Fontaine a
construit un anneau $B_{\mathrm{cris}}$ tel que $D_{\mathrm{cris}}(M) = (M
\otimes B_{\mathrm{cris}})^{\pi}$ soit, sur les représentations dites
\emph{cristallines}, un $\otimes$-foncteur exact à valeurs dans les
$\varphi$-modules filtrés, avec $\mathrm{gr}\, D_{\mathrm{dR}} \simeq
D_{\mathrm{HT}}$ ; les $V_p(G)$ sont cristallines, et
$D_{\mathrm{cris}}(V_p(G))$ redonne le module de Dieudonné.\footnote{Le
rapprochement entre cette Question et le « foncteur mystérieux » est le nôtre :
aucun feuillet du dossier n'emploie l'expression. Les constructions de
Fontaine datent de 1979 à 1982 ; les pages n'y renvoient pas. Rien dans le
dossier ne dit si les feuillets sont d'avant ou d'après 1970.}

\pagerange{24}{26}
\section*{III. La liste des données (pages 24 à 26)}

Les trois pages au crayon changent de point de vue. Il ne s'agit plus de la
catégorie $\mathcal{M}_K$ et de son groupe, mais d'un groupe algébrique $G$
\emph{quelconque} sur $\mathbb{Q}_p$ : que faut-il se donner pour avoir, à
valeurs dans $G$, l'équivalent de tout ce que porte un groupe de
Barsotti-Tate ? La réponse est une liste numérotée, commencée deux fois sur
trois pages. On rappelle que, sur ces pages, la lettre $L$ de la page désigne notre $K$, et
sa lettre $K$ notre $K_0$.

\pagerange{24}{24}
\subsection*{Première rédaction (page 24)}

Les données sont : 1) un groupe algébrique $G$ sur $\mathbb{Q}_p$ ; 2) un
homomorphisme continu $\rho : \pi \to G(\mathbb{Q}_p)$, d'où, comme à la
page 22, le groupe $G^{!}_{\Omega}$, le torseur $\mathfrak{P}$ et le groupe
$G_{\Omega}$ ; 3) un cocaractère $j^{!} : \mathbb{G}_{m,\Omega} \to
G^{!}_{\Omega}$, d'où $Q_{\Omega}$, $G^{\mathrm{Hdg}}_{\Omega}$ et
$j_{\Omega}$ ; 4) un torseur $Q_K$ sous $G_K$ qui descend $Q_{\Omega}$, d'où le
groupe $G^{\mathrm{Hdg}}_K$, un sous-groupe $U^{\mathrm{Hdg}}_K$ et le
cocaractère $j_K$. Une accolade embrasse 3) et 4) et porte en marge le mot
« unique ».\footnote{Le « 4) » est d'abord écrit puis biffé, et réécrit à la
ligne suivante ; le $\mathcal{T}$ de 3) est écrit par-dessus une autre
lettre, et l'indication « $\mathrm{int}$ » qui précise que $G^{!}_{\Omega}$ est
tordu par automorphismes intérieurs est une lecture incertaine.}

La cinquième donnée n'aboutit pas, et ce qu'elle cherche est intéressant. Elle
fait apparaître, avant toute définition, un groupe $\Gamma$ contenant $\pi$,
et considère un ensemble $\Delta$ d'opérations de $\Gamma$ sur
$(\Omega, Q_K, G_K)$ qui prolongent celle de $\pi$ et sont compatibles à
l'action de $\Gamma$ sur $\Omega$. Ces opérations ne sont pas uniques : elles
forment une orbite sous le groupe $U^{\mathrm{Hdg}}_K(K)$, qui agit par
conjugaison,
\[
(u \cdot \rho)(\gamma) = \tau(u)\, \rho(\gamma)\, \tau(u)^{-1}.
\]
La version finale de la donnée, non biffée, demande un torseur à droite
$\Delta$ sous $U^{\mathrm{Hdg}}_K(K)$ et une application équivariante de
$\Delta$ dans l'ensemble de ces opérations.\footnote{La première ligne de 5)
est biffée d'un trait horizontal, et le début de la deuxième ne se lit pas ;
tout le passage qui va de « tel que » à la formule de conjugaison est en outre
annulé par trois longs traits obliques, et un « NB » qui suit est biffé
lui-même. La lettre $\tau$ n'est définie nulle part : c'est vraisemblablement
l'action de $U^{\mathrm{Hdg}}$ par automorphismes. Dans la phrase finale,
« à droite $\Delta$ » est ajouté en interligne au-dessus d'un $R_0$ qu'il
remplace peut-être, précédé d'un « commutatif » incertain. La même phrase
définit l'action à droite par « $u.\alpha = \tau(\alpha)^{-1}\, u\,
\tau(\alpha)$ » : les lettres y sont manifestement échangées, et l'action
cohérente avec la formule précédente est $\alpha \cdot u =
\tau(u)^{-1}\, \alpha\, \tau(u)$.} L'idée est là : les façons d'ajouter
à l'action de $\pi$ l'action d'un groupe plus gros ne forment pas un point,
mais une orbite. La page 24 le \emph{postule} ; les pages 27 et 28
demandent si c'est vrai.

\pagerange{25}{25}
\subsection*{Seconde rédaction : le groupe, la représentation, la graduation (page 25)}

La page recommence à 1). Le groupe $G$ sur $\mathbb{Q}_p$ « dépend en fait du
choix de $\Omega$ », et un renvoi annonce l'explication, qui vient en bas de
page. Avec 2), l'homomorphisme $\rho$, le groupe $\pi = \mathrm{Aut}_K(\Omega)$
agit sur $G$ par $\gamma \mapsto \mathrm{int}(\rho(\gamma))$, et sur
$G_{\Omega}$ \emph{diagonalement} : c'est $G^{!}_{\Omega}$.\footnote{La page
écrit $\pi = \mathrm{Aut}_L(\Omega)$ : c'est le premier des indices de
l'échange de $K$ et $L$.} La donnée 3) est le cocaractère $j^{!} :
\mathbb{G}_{m,\Omega} \to G^{!}_{\Omega}$, « invariant par $\pi$ ». On pose
$\mathcal{T}$, le $\mathbb{G}_m$-torseur associé au module inversible
$\Omega(1)$, puis $Q_{\Omega} = \mathcal{T} \wedge \mathfrak{P}$, torseur à
droite sous $G_{\Omega}$, et $G^{\mathrm{Hdg}}_{\Omega}$ son groupe
d'automorphismes, qui correspond aux foncteurs fibres gradués sur les
représentations de $G$ ; le foncteur fibre de $Q_{\Omega}$ est
\[
M \longmapsto \bigoplus_i \mathrm{gr}^{i}_{j}\bigl(\omega^{!}(M)\bigr)
\otimes_{\Omega} \Omega(-i) = \mathrm{Hdg}_{\Omega}(M),
\]
la formule même de la section II.\footnote{Une ligne entière, entre « Posons »
et « torseur », est raturée jusqu'à l'illisible ; on y devine
$\mathfrak{P}$, $Q$ et $\mathbb{G}_{m,\Omega}$. La formule $Q_{\Omega} =
\mathcal{T} \wedge \mathfrak{P}$ est écrite au-dessus, dans un cadre, et
« Posons » est en interligne. Devant « Hdg », un trait vertical peut être un
$T$ ou le crochet fermant de la ligne précédente. En marge, un $G_{\Omega}$
incertain.}

\emph{Axiome.} Le torseur $Q_{\Omega}$ descend à $K$ en un $G_K$-torseur $Q_K$,
unique à isomorphisme unique près, et de même $G^{\mathrm{Hdg}}_{\Omega}$ et
$j_{\Omega}$ en $G^{\mathrm{Hdg}}_K$ et $j_K : \mathbb{G}_{m,K} \to
G^{\mathrm{Hdg}}_K$ ; cela revient à se donner le foncteur fibre gradué
$\mathrm{Hdg}_{j_K}$ sur $K$.\footnote{La page écrit « se descend à $L$ ». Au
début de l'Axiome, le mot « torseur » est biffé sous une surcharge entre
parenthèses non lue ; $Q_L$ et $G_L$ sont repris dans un ovale.}

Une remarque de marge précise : \emph{un $j$ satisfaisant aux conditions est
unique, et son existence signifie que les $M_{\Omega}$ sont de Hodge}.
C'est exact, à condition d'entendre par « conditions » la donnée 3)
\emph{et} l'Axiome.\footnote{La marge se termine par un mot illisible.} Un
cocaractère invariant par $\pi$ ne suffit pas à lui seul : le cocaractère
trivial l'est toujours. C'est l'Axiome --- la trivialité de chaque
$\mathrm{gr}^i \otimes \Omega(-i)$ --- qui force $\omega^{!}(M)_i \simeq
\Omega(i)^{h_i}$, c'est-à-dire la propriété de Hodge-Tate, et l'unicité vient
alors de $\mathrm{Hom}_{\pi}(\Omega(i), \Omega(i')) = 0$ pour $i \neq i'$.

Le bas de la page donne l'explication promise, marquée d'un « 2° » cerclé. Il
y a deux foncteurs fibres des représentations de $G$ dans les $\Omega$-espaces
à action semi-linéaire de $\pi$ : a) $M \mapsto \omega(M)$, où $\pi$ n'agit que
sur $\Omega$, et b) $M \mapsto \omega^{!}(M)$, où il agit diagonalement. On a
donc un torseur à droite sous $G_{\Omega}$, qui est $\mathfrak{P}$, donné par
le $1$-cocycle $\rho : \pi \to G(\mathbb{Q}_p) \subset G(\Omega)$, et son
groupe d'automorphismes est $G^{!}_{\Omega}$. C'est ce qu'on a utilisé à la
section II pour lire le tableau de la page 22.\footnote{« 2° » est lu d'après
le renvoi « cf plus bas 2° » du haut de la page ; en a), « par transport » est
une lecture incertaine, de même que l'indice $(\Omega, \pi)$ qui suit
$\mathrm{Aut}$ dans la dernière phrase.}

\pagerange{26}{26}
\subsection*{Filtration, descente, Frobenius (page 26)}

La page 26 continue la page 25 à partir de 4).

\emph{4) La filtration.} Soit $U^{\mathrm{Hdg}}_K \subset G^{\mathrm{Hdg}}_K$ le
sous-groupe \emph{défini par} $j_K$ : le radical unipotent du parabolique
associé, c'est-à-dire le sous-groupe des automorphismes qui respectent la
filtration définie par la graduation et induisent l'identité sur le
gradué.\footnote{La page dit seulement « défini par $j_L$ » ; la description
est la nôtre, et c'est la seule qui rende correct ce qui suit.} On a un
torseur à droite $R_K$ sous $U^{\mathrm{Hdg}}_K$, « canonique (mais
trivial) », d'où un foncteur fibre \emph{filtré} sur $K$ :
\[
\mathrm{DR}(M) = \bigl(R_K \wedge^{U^{\mathrm{Hdg}}_K} Q_K\bigr)
\wedge^{G_K} (T_p(M) \otimes K).
\]
Voici ce qui justifie la formule. Tordre le foncteur gradué
$T_{\mathrm{Hdg}}$, vu comme filtré, par un $U^{\mathrm{Hdg}}_K$-torseur donne
un foncteur fibre filtré dont le gradué est encore $T_{\mathrm{Hdg}}$, et
inversement tout foncteur fibre filtré sur $K$ muni d'un isomorphisme
$\mathrm{gr} \simeq T_{\mathrm{Hdg}}$ s'obtient ainsi : $R_K$ est le torseur
de ses scindages compatibles au produit tensoriel. En caractéristique nulle,
un tel scindage existe toujours, et $U^{\mathrm{Hdg}}_K$ étant unipotent, tout
torseur sous lui est trivial ; d'où le « mais trivial ». Le foncteur filtré
est isomorphe au foncteur scindé, mais pas canoniquement.\footnote{« canonique »
est une lecture incertaine, et l'identification de $R_K$ au torseur des
scindages est la nôtre. L'existence des scindages des foncteurs fibres filtrés
en caractéristique nulle est établie par Saavedra (1972). Sur la page, une
première écriture $T_p(M)_L$ est biffée après le signe $=$.}

Le $G_K$-torseur qui correspond au foncteur $\mathrm{DR}$ est
\[
P_K = R_K \wedge^{U^{\mathrm{Hdg}}_K} Q_K ,
\]
et l'on pose $G'_K = \mathrm{Aut}(P_K)$, forme intérieure de $G_K$.\footnote{La
page note ici le groupe contracté $U'_L$, avec un seul accent, là où partout
ailleurs elle écrit $U''_L$ ; c'est le même groupe.} La page ajoute
« $G'_K \simeq G_K$ non canoniquement » ; cela suppose que le torseur $P_K$
soit trivial, ce qui a lieu par exemple pour $G = \mathrm{GL}_n$ (théorème 90
de Hilbert), mais pas pour tout $G$.\footnote{Hypothèse ajoutée par nous : pour
un groupe dont $H^1(K, G)$ n'est pas nul, un orthogonal par exemple, rien dans
les données n'oblige $P_K$ à être trivial.}

\emph{5) La descente.} Une structure de descente de $K$ à $K_0$ sur le
torseur $P_K$, c'est-à-dire un $G_{K_0}$-torseur $P_{K_0}$ dont $P_K$ est
l'extension des scalaires ; d'où $G'_{K_0}$ et $P_{K_0}$.\footnote{La page
écrit « Structure de descente de $L : K$ sur le torseur $P_L$, d'où $G''$,
$P_K$, $G_K$ ». Un « 5) » placé une ligne plus haut est biffé.}

\emph{6) Le Frobenius.} Un isomorphisme de $G_{K_0}$-torseurs donné
\[
P_{K_0} \simeq P_{K_0}^{(\sigma)},
\]
où $P^{(\sigma)}$ est le torseur déduit de $P$ par le changement de base
$\sigma : K_0 \to K_0$.\footnote{L'exposant est un signe entre parenthèses
qui tient du $\sharp$ et du $H$, et la transcription ne le lit pas avec
certitude. La lecture $(\sigma)$ est la nôtre : c'est la seule qui fasse de 6)
une structure de Frobenius, ce que demande la donnée b) de la page 23.}

Les données 5) et 6) réunies disent que le foncteur fibre $\mathrm{DR}$ provient
d'un foncteur fibre sur $K_0$ muni d'un Frobenius : c'est la version à
valeurs dans $G$ des données a) à c) de la page 23. On les appelle aujourd'hui
des \emph{isocristaux à $G$-structure}, et le couple formé du Frobenius et du
cocaractère de la filtration est ce qu'on étudie sous le nom de couple
$(b, \mu)$.\footnote{Les isocristaux à $G$-structure sont dus à R. Kottwitz
(1985) ; le couple $(b, \mu)$ est au centre des travaux de Rapoport et Zink
dans les années 1990. Les pages ne citent rien de cela, et on ne les
présente pas comme y conduisant.}

\emph{L'extension $\Gamma$.} Pour finir, la page propose de formuler 5) et 6)
ensemble. On considère le torseur $P_{\Omega}$ déduit de $P_K$, sur lequel
$\pi$ agit par son action sur $\Omega$, et « l'extension $\Gamma$ de
$\mathbb{Z}$ par $\tilde{\pi}$ », qui contient $\pi$, et l'on déclare que
$\Gamma$ \emph{opère} sur $P_{\Omega}$ en prolongeant l'opération de
$\pi$.\footnote{La page note $\tilde{\pi} = \mathrm{Gal}(\bar{K}/\tilde{K})$,
lecture incertaine ; « extension de $\mathbb{Z}$ par $\tilde{\pi}$ » est à
prendre au sens de Grothendieck, $1 \to \tilde{\pi} \to \Gamma \to \mathbb{Z}
\to 1$. La formule qui précède, $P_{\Omega} = R_L \wedge^{U''_{\Omega}}
Q_{\Omega}$, est surchargée, et son second membre s'arrête après le signe
$\wedge$.} La page ne définit pas $\Gamma$ davantage. Une réalisation qui rend
l'énoncé exact est la suivante : $\Gamma$ est le groupe des couples
$(\gamma, n)$ formés d'un automorphisme continu $\gamma$ de $\Omega$ et d'un
entier $n$ tels que $\gamma$ induise $\sigma^n$ sur $K_0$. On a alors
$1 \to \pi' \to \Gamma \to \mathbb{Z} \to 1$, avec $\pi' =
\mathrm{Gal}(\bar{K}/K_0) \supset \pi$, le groupe que la page 22 avait posé
sans s'en servir. Prolonger à $\pi'$ l'action de $\pi$ sur $P_{\Omega}$, c'est
la descente 5), puisque $\Omega^{\pi'} = K_0$ ; faire agir les éléments qui
relèvent $\sigma$, c'est le Frobenius 6).\footnote{Cette réalisation de
$\Gamma$, et l'identification de $\tilde{\pi}$ au $\pi'$ de la page 22, sont
nôtres. Lorsque $k$ est fini, $\sigma$ est d'ordre fini sur $K_0$, et c'est
pourquoi l'on prend des couples plutôt que les seuls automorphismes.} L'idée de
coder dans un seul groupe l'action de Galois et celle du Frobenius est celle
qu'on retrouve, sous d'autres formes, dans le groupe de Weil et dans les
anneaux de Fontaine, où agissent ensemble $\pi$ et un Frobenius.

\pagerange{27}{28}
\section*{IV. L'unicité : une question de cohomologie non abélienne (pages 27 et 28)}

Les deux dernières pages au crayon traduisent la page 26 en cohomologie
galoisienne non commutative, et posent la question de l'unicité.

\subsection*{L'extension et ses sections}

Soit $A = G'(\Omega)$ le groupe des automorphismes du $G_{\Omega}$-torseur
$P_{\Omega}$.\footnote{La page écrit $G'_S(S)$, avec un $S$ cursif qui n'est
défini nulle part dans le dossier. L'égalité $G'_S(S)^{\pi} = G'_L(L)$ qu'elle
écrit plus loin exige que les invariants de $\pi$ sur $S$ soient le corps de
base ; c'est le cas de $\Omega$, et nous écrivons $\Omega$.} Soit
$\mathcal{E}$ le groupe des couples $(\gamma, u)$ formés d'un $\gamma \in
\Gamma$ et d'un automorphisme $u$ de $P_{\Omega}$, compatible à l'action de
$G_{\Omega}$ et semi-linéaire relativement à l'action de $\gamma$ sur $\Omega$.
Comme $\Omega$ est algébriquement clos, $P_{\Omega}$ est trivial, et tout
$\gamma$ se relève ; on a donc une extension
\[
1 \longrightarrow A \longrightarrow \mathcal{E} \longrightarrow \Gamma
\longrightarrow 1 . \qquad (\mathcal{E})
\]
Ses sections --- les homomorphismes $s : \Gamma \to \mathcal{E}$ qui
relèvent l'identité --- sont exactement les actions de $\Gamma$ sur
$P_{\Omega}$ compatibles à la structure de torseur. Une section $s$ fait
opérer $\Gamma$ sur $A$ par conjugaison, et cette opération provient,
remarque la page, d'une opération de $\Gamma$ sur le schéma en groupes
$G'_{\Omega}$ lui-même. Les autres sections s'écrivent $\varphi \cdot s$,
où $\varphi$ parcourt les $1$-cocycles de $\Gamma$ à valeurs dans $A$ (pour
cette opération), et $A$ opère sur l'ensemble des sections par conjugaison,
sans restriction : deux sections sont conjuguées si et seulement si le cocycle
qui fait passer de l'une à l'autre est cohomologue à $1$. Les classes de
conjugaison de sections forment donc l'ensemble pointé $H^1(\Gamma,
A)$.\footnote{Une phrase commencée, « Supposons que $\pi \subset \Gamma$
opère sur », est biffée en tête de la page 27. « est coh. à 1 » est ajouté en
marge avec un trait de renvoi, et « aucune restriction » est une lecture
incertaine.}

Soit maintenant $\alpha_0 : \pi \to \mathcal{E}$ une section partielle,
au-dessus de $\pi$ --- l'action donnée de $\pi$ sur $P_{\Omega}$. Le
stabilisateur de $\alpha_0$ dans $A$ est le groupe des invariants $A^{\pi}$,
pour l'opération définie par $\alpha_0$, et
\[
A^{\pi} = G'(\Omega)^{\pi} = G'(K),
\]
toujours parce que $\Omega^{\pi} = K$.\footnote{La page écrit $G'_S(S)^{\pi} =
G'_L(L)$. Le passage qui introduit ce groupe contient un mot illisible et deux
lectures incertaines, « ens. » et « comme on pense ».} Ce groupe opère sur
l'ensemble des sections $\alpha$ qui prolongent $\alpha_0$. Si l'on en a une,
les autres s'en déduisent par les $1$-cocycles \emph{triviaux sur} $\pi$ :
\[
\begin{cases}
\varphi(\gamma\gamma') = \varphi(\gamma)\; {}^{\gamma}\varphi(\gamma'), \\
\varphi(\gamma) = e \quad \text{pour } \gamma \in \pi .
\end{cases}
\]

\subsection*{La Question}

\emph{Question} (page 28). Les $\alpha$ qui prolongent $\alpha_0$ sont-ils
conjugués sous $A$ --- ce qui revient au même que sous $A^{\pi}$, puisqu'un
élément qui conjugue deux prolongements de $\alpha_0$ centralise
$\alpha_0(\pi)$ ? La page observe qu'il suffirait de savoir que, pour une
telle section $\alpha$ fixée, l'application
\[
H^1(\Gamma, A) \longrightarrow H^1(\pi, A)
\]
est « injective », au sens des ensembles pointés : de noyau réduit à
l'élément neutre. Les solutions formeraient alors un espace homogène sous
$A^{\pi} = G'(K)$.\footnote{Trois traits verticaux en marge de gauche
marquent la Question. Un long trait vertical traverse la page de haut en bas,
longé d'un zigzag dans sa moitié inférieure ; ils séparent plutôt qu'ils ne
biffent.}

La condition est en fait nécessaire et suffisante, dès qu'un prolongement
$\alpha$ existe. Les éléments du noyau sont représentés par les cocycles
triviaux sur $\pi$ --- un cocycle dont la restriction à $\pi$ est un cobord se
modifie par ce cobord --- ; et deux tels cocycles sont cohomologues sur
$\Gamma$ si et seulement s'ils le sont par un élément de $A^{\pi}$. Le noyau
est donc en bijection avec l'ensemble des orbites de $A^{\pi}$ sur les
prolongements de $\alpha_0$, et il est trivial si et seulement si ces
prolongements forment une seule orbite.\footnote{La page dit « il suffirait » ;
la réciproque est notre ajout.}

\subsection*{Ce que l'on peut répondre}

Telle qu'elle est posée, avec les seules données 1) à 6), la réponse est non.
Prenons $K = K_0 = \mathbb{Q}_p$, $G = \mathbb{G}_m$, $\rho$ trivial et $j$
trivial. Toutes les données jusqu'à 5) sont alors triviales, $P_{\Omega}$ est
le torseur trivial, $A = \Omega^{*}$ et $A^{\pi} = \mathbb{Q}_p^{*}$. Comme
$\sigma$ est l'identité de $\mathbb{Q}_p$, le groupe $\Gamma$ est le produit
$\pi \times \mathbb{Z}$, et un prolongement de $\alpha_0$ est déterminé par
l'image du générateur de $\mathbb{Z}$, qui doit commuter à $\pi$ : c'est la
multiplication par un élément $b$ de $\Omega^{*}$ invariant par $\pi$, donc de
$\mathbb{Q}_p^{*}$. C'est exactement la donnée 6), un Frobenius $b$ sur le
$\mathbb{Q}_p$-espace de dimension 1. Et $A^{\pi}$, commutatif, agit
trivialement : chaque $b \in \mathbb{Q}_p^{*}$ est une orbite à lui seul. Le
noyau de $H^1(\Gamma, A) \to H^1(\pi, A)$ est $\mathbb{Q}_p^{*}$ tout
entier.\footnote{L'exemple est le nôtre, de même que la réalisation de
$\Gamma$ qu'il utilise.}

La raison est dans la liste elle-même. La représentation $\rho$ n'y intervient
qu'à travers $\mathfrak{P}$ et la graduation $j$, c'est-à-dire à travers sa
décomposition de Hodge-Tate ; rien n'y relie le Frobenius à $\rho$. Dans
l'exemple, $b = p$ donne l'isocristal de pente $1$, qui ne peut pas
correspondre à une représentation de poids $0$ ; mais même la condition de
Fontaine qui exclut ce cas --- l'admissibilité faible, qui demande ici que $b$
soit une unité --- laisse subsister tous les $b \in \mathbb{Z}_p^{*}$, qui
sont les isocristaux des caractères non ramifiés : d'autres représentations,
que les données de Hodge-Tate ne distinguent pas de la représentation
triviale. Ce qui manque à la liste est un lien direct entre $\rho$ et le
Frobenius, et c'est ce que fournit l'anneau $B_{\mathrm{cris}}$ de Fontaine,
où agissent à la fois $\pi$ et un Frobenius, et qui fait de
$D_{\mathrm{cris}}$ une fonction de la représentation.\footnote{Sur un corps
résiduel algébriquement clos, en revanche, les unités de $K_0$ sont toutes de
la forme $\sigma(a)/a$, et dans l'exemple de rang 1 l'admissibilité faible
suffit à l'unicité. La question de la page est donc sensible au corps
résiduel, ce qu'elle ne dit pas.}

\pagerange{28}{28}
\section*{V. Tores sur un corps $p$-adique (bas de la page 28)}

Sous un trait horizontal et une lettre H au crayon, la moitié inférieure de la
page 28 est un calcul à l'encre, indépendant de ce qui précède. Il porte sur un
corps $K$ fini sur $\mathbb{Q}_p$ --- l'hypothèse est nécessaire pour que le
tore suivant soit défini --- et sur le tore
\[
T_K = \mathrm{R}_{K/\mathbb{Q}_p}\, \mathbb{G}_m, \qquad
T_K(\mathbb{Q}_p) \simeq K^{*},
\]
restriction des scalaires à la Weil de $\mathbb{G}_m$, que la page écrit
$\prod_{K/\mathbb{Q}_p} \mathbb{G}_m$.

La marge nomme l'objet du calcul : le \emph{modèle de Néron} de ce tore,
qu'il faut entendre localement de type fini, puisque $K^{*}$ n'est pas
borné ; on le note $N$, sur $\mathbb{Z}_p$. Sa
composante neutre est $N^{\circ} = \mathrm{R}_{\mathcal{O}_K/\mathbb{Z}_p}\,
\mathbb{G}_m$, et l'on a une suite exacte
\[
0 \longrightarrow N^{\circ} \longrightarrow N \longrightarrow \Phi
\longrightarrow 0,
\]
où le groupe des composantes $\Phi$ est un groupe étale concentré sur la fibre
spéciale. La propriété de Néron donne $N(\mathbb{Z}_p) = T_K(\mathbb{Q}_p) =
K^{*}$, et la suite des points entiers est la suite de la valuation,
\[
1 \longrightarrow \mathcal{O}_K^{*} \longrightarrow K^{*}
\xrightarrow{\ v\ } \mathbb{Z} \longrightarrow 0,
\]
que la page écrit en colonne.\footnote{La page note $\sigma$ l'anneau des
entiers, en indice : $\prod_{\sigma/\mathbb{Z}_p}\mathbb{G}_m$,
$g_{\sigma}$, $\sigma^{*}_K$ ; nous écrivons $\mathcal{O}_K$, la lettre
$\sigma$ étant ici le Frobenius. Elle note $g_K$ le tore et $g_{\sigma}$ la
composante neutre de son modèle. Dans le tableau de gauche, un signe
ressemblant à « m », non lu, précède la troisième et la quatrième ligne ; la
dernière ligne est étiquetée « $\otimes_{\sigma} K$ ». Ces marques ne sont pas
utilisées ici.}

La page écrit $\mathbb{Z}$ pour le groupe des composantes. C'est exact lorsque
$K/\mathbb{Q}_p$ est totalement ramifiée ; en général, $\Phi$ est la
restriction des scalaires de $\mathbb{Z}$ du corps résiduel de $K$ à
$\mathbb{F}_p$, dont les points sur $\mathbb{F}_p$ forment encore
$\mathbb{Z}$.\footnote{Précision nôtre ; la page ne dit rien du degré
résiduel.} Plaçons-nous dans le cas totalement ramifié. Le choix d'une
uniformisante $\varpi$ de $K$ fournit un morphisme
\[
\tilde{N} = N^{\circ} \times \mathbb{Z} \longrightarrow N, \qquad (u, n)
\longmapsto u\,\varpi^{n},
\]
où $\mathbb{Z}$ est le groupe constant sur $\mathbb{Z}_p$. C'est un
isomorphisme sur les points entiers, $\mathcal{O}_K^{*} \times \mathbb{Z}
\simeq K^{*}$, et sur les fibres spéciales, mais \emph{pas} sur les fibres
génériques : celle de $\tilde{N}$ est $T_K \times \mathbb{Z}$, celle de $N$
est $T_K$. La page écrit précisément cette fibre générique,
$\tilde{N}_K \simeq g_K \times \mathbb{Z}_K$, dont les points forment
$K^{*} \times \mathbb{Z}$, et l'inclusion $\Gamma(\tilde{N}) \hookrightarrow
\Gamma(\tilde{N}_K)$, c'est-à-dire $\mathcal{O}_K^{*} \times \mathbb{Z}
\hookrightarrow K^{*} \times \mathbb{Z}$. Les deux lignes centrales du tableau
de la page se lisent ainsi :\footnote{Le tableau de la page a quatre lignes et
une cinquième amorcée : au-dessus, $\prod_{\sigma/\mathbb{Z}_p}\mathbb{G}_m
\to T_{\sigma_K} \to \mathbb{Z} \to 0$, qui redit la ligne de $N$ ; en
dessous, la fibre générique $0 \to g_K \to \tilde{N}_K \to \mathbb{Z}_K \to
0$ et l'isomorphisme $\tilde{N}_K \simeq g_K \times \mathbb{Z}_K$. On ne
redessine que les deux lignes qui portent le calcul. La page trace entre les
deux $g_{\sigma}$ une double flèche ; c'est une égalité, rendue ici par une
liaison sans tête.}

\begin{tikzcd}
0 \arrow[r] & N^{\circ} \arrow[r] \arrow[d, no head] & \tilde{N} \arrow[r] \arrow[d] & \mathbb{Z} \arrow[r] \arrow[d] & 0 \\
0 \arrow[r] & N^{\circ} \arrow[r] & N \arrow[r] & \Phi \arrow[r] & 0
\end{tikzcd}

La flèche oblique que la page tire de $K^{*}$ vers $\Gamma(\tilde{N}_K)$ est
marquée $\pi$. Ce $\pi$ n'est pas le groupe de Galois des pages précédentes :
c'est vraisemblablement l'uniformisante, dont dépend la flèche $x \mapsto
(x\,\varpi^{-v(x)}, v(x))$.\footnote{La transcription propose déjà cette
lecture. La flèche montante de $\Gamma(\tilde{N})$ vers $\Gamma(N)$ porte un
« 2 » de lecture douteuse, qu'on ne reprend pas : la flèche est un
isomorphisme. À droite de la colonne, un second $\mathbb{Z}$ flanqué de
flèches montantes n'a pas de source écrite.}

Le lien de ce calcul avec le reste du dossier n'est pas écrit. Le seul qu'on
puisse proposer est que $T_K$ est le tore à travers lequel Serre, dans son
livre sur les représentations $\ell$-adiques abéliennes (1968), décrit par le
corps de classes local les représentations abéliennes localement algébriques
de $\pi$ ; la page ne le dit pas.\footnote{Le rapprochement est le nôtre.}

\end{document}
